\DocumentMetadata{
  pdfversion=2.0,
  pdfstandard=ua-2,
  lang=en-US,
  tagging=on,
  tagging-setup={math/setup=mathml-SE,tabsorder=structure}
}
\documentclass[11pt,letterpaper]{article}
\usepackage[margin=0.68in,includefoot]{geometry}
\usepackage{newcomputermodern}
\usepackage{amsmath,amscd,mathtools}
\usepackage{tikz,adjustbox}
\providecommand{\square}{\mdlgwhtsquare}
\usepackage[hidelinks]{hyperref}
\hypersetup{
  pdftitle={Analysis PhD Exam, January 2024},
  pdfauthor={University of Florida Department of Mathematics},
  pdfsubject={Graduate mathematics examination},
  pdfdisplaydoctitle=true,
  bookmarksopen=true
}
\setcounter{secnumdepth}{0}
\setlength{\parindent}{0pt}
\setlength{\parskip}{0.28em}
\setlength{\emergencystretch}{3em}
\setlength{\leftmargini}{2.15em}
\setlength{\leftmarginii}{2.65em}
\setlength{\leftmarginiii}{3em}
\pagestyle{empty}
\raggedbottom
\allowdisplaybreaks
\NewTaggingSocketPlug{block/list/label}{examproblem}{%
  \tagstructbegin{tag=Lbl}\tagstructend%
  \tagstructbegin{tag=\LBody}%
  \tagstructbegin{tag=\ProblemHeadingTag,title-o={\ProblemHeadingText}}%
  \tagmcbegin{tag=\ProblemHeadingTag,actualtext={\ProblemHeadingText}}%
  #2%
  \tagmcend%
  \tagstructend%
}
\newenvironment{examproblems}[2]{%
  \def\ProblemHeadingTag{#2}%
  \AssignTaggingSocketPlug{block/list/label}{examproblem}%
  \begin{enumerate}%
  \setcounter{enumi}{#1}%
  \renewcommand{\labelenumi}{\textbf{\theenumi.}}%
  \setlength{\topsep}{0.35em}%
  \setlength{\partopsep}{0pt}%
  \setlength{\itemsep}{0.48em}%
  \setlength{\parsep}{0.18em}%
}{\end{enumerate}}
\newenvironment{examsubparts}{%
  \AssignTaggingSocketPlug{block/list/label}{default}%
  \begin{enumerate}%
  \setlength{\topsep}{0.18em}%
  \setlength{\partopsep}{0pt}%
  \setlength{\itemsep}{0.16em}%
  \setlength{\parsep}{0.1em}%
}{\end{enumerate}}
\newenvironment{examinstructions}{%
  \par\begingroup\itshape
}{%
  \par\endgroup\vspace{0.18em}
}
\makeatletter
\renewcommand\subsection{\@startsection{subsection}{2}{0pt}%
  {-1.25ex plus -.25ex minus -.1ex}%
  {0.55ex plus .1ex}%
  {\normalfont\large\bfseries}}
\renewcommand{\@maketitle}{%
  \newpage\null\vskip -1em%
  {\footnotesize\bfseries
    \href{https://www.ufl.edu/}{University of Florida}, \href{https://math.ufl.edu/}{Department of Mathematics}\par}%
  \vskip -0.5em%
  \tagstructbegin{tag=H1,title={Analysis PhD Exam, January 2024}}%
  \tagpdfparaOff%
  \tagmcbegin{tag=H1}%
    {\LARGE\bfseries\href{https://gma.math.ufl.edu/exams/analysis-phd/}{Analysis PhD Exam}, January 2024\par}%
  \tagmcend%
  \tagpdfparaOn%
  \tagstructend%
  \vskip 0.25em%
  \tagmcbegin{artifact}%
    \hrule height 1pt%
  \tagmcend%
  \vskip 1em%
}
\makeatother
\title{Analysis PhD Exam, January 2024}
\author{}
\date{}
\begin{document}
\maketitle
\thispagestyle{empty}
\begin{examinstructions}
Write solutions in a neat and logical fashion, giving complete reasons for all steps. Attempt SIX problems.
\end{examinstructions}
\begin{examproblems}{0}{H2}
\def\ProblemHeadingText{Problem 1.}
\item
Let \((X,\mathcal{M},\mu)\) be a finite measure space. Let \((f_n)\) be a sequence of measurable functions and suppose that \(\sum_n \int |f_n|\,d\mu<\infty\). Consider the series \(\sum_n f_n(x)\). Does this series converge… Prove your claims.
\begin{examsubparts}
\item[\textnormal{a)}]
pointwise a.e.?
\item[\textnormal{b)}]
in measure?
\item[\textnormal{c)}]
in the \(L^1\) norm?
\end{examsubparts}
\def\ProblemHeadingText{Problem 2.}
\item
State Tonelli’s theorem, and give an example to show that the~\(\sigma\)-finiteness hypothesis is necessary.
\def\ProblemHeadingText{Problem 3.}
\item
This problem has two parts.
\begin{examsubparts}
\item[\textnormal{a)}]
Let~\(\nu\) be a signed measure on \((X,\mathcal{M})\). Define what it means for a set to be positive, negative or null for~\(\nu\), and state the Hahn decomposition theorem.
\item[\textnormal{b)}]
Prove that if~\(\nu\) is a signed measure, then there exist unique positive measures \(\nu^+,\nu^-\) so that \(\nu=\nu^+-\nu^-\) and \(\nu^+\perp\nu^-\).
\end{examsubparts}
\def\ProblemHeadingText{Problem 4.}
\item
Let \(\varphi_n:\mathbb{R}\to\mathbb{C}\) be a sequence of Lebesgue measurable functions and suppose that there is a number~\(M\) such that \(\|\varphi_n\|_\infty\le M\) for all~\(n\). Prove that if
\[
\int_a^b \varphi_n(t)\,dt\to 0
\]
 for every interval \([a,b]\subset\mathbb{R}\), then
\[
\int_{-\infty}^{\infty}\varphi_n(t)f(t)\,dt\to 0
\]
 for every \(f\in L^1(\mathbb{R})\).
\def\ProblemHeadingText{Problem 5.}
\item
Let~\(\mu\) and~\(\nu\) be finite, positive measures defined on the same measurable space \((X,\mathcal{M})\), and suppose that \(\nu\ll\mu\).
\begin{examsubparts}
\item[\textnormal{a)}]
Prove that if \(L^1(\mu)\subset L^1(\nu)\), then the inclusion map \(\iota:L^1(\mu)\hookrightarrow L^1(\nu)\) is necessarily bounded.
\item[\textnormal{b)}]
Prove that if \(L^1(\mu)\subset L^1(\nu)\) then \(\frac{d\nu}{d\mu}\in L^\infty(\mu)\).
\end{examsubparts}
\def\ProblemHeadingText{Problem 6.}
\item
Let \(\mathcal{X},\mathcal{Y}\) be Banach spaces and \(T:\mathcal{X}\to\mathcal{Y}\) a linear map. Prove that if \(f\circ T\in\mathcal{X}^*\) for all \(f\in\mathcal{Y}^*\), then~\(T\) is bounded.
\def\ProblemHeadingText{Problem 7.}
\item
Short answer.
\begin{examsubparts}
\item[\textnormal{a)}]
Explain why there is no norm on \(c_{00}\) (the space of finitely nonzero sequences) which makes it into a Banach space.
\item[\textnormal{b)}]
Explain why there is no bounded, surjective linear map from \(c_0\) onto \(\ell^\infty\).
\item[\textnormal{c)}]
Explain why the norm in \(L^1(\mathbb{R})\) cannot be given by an inner product.
\end{examsubparts}
\def\ProblemHeadingText{Problem 8.}
\item
Let~\(H\) be a Hilbert space with orthonormal basis \(\{e_k\}_{k=1}^\infty\). Let \((h_n)\) be a sequence in~\(H\) and let \(h\in H\). Prove that the following statements are equivalent:
\begin{examsubparts}
\item[\textnormal{1)}]
\(\langle h_n,g\rangle\to\langle h,g\rangle\) for all \(g\in H\).
\item[\textnormal{2)}]
\(\langle h_n,e_k\rangle\to\langle h,e_k\rangle\) for all \(k=1,2,\ldots\), AND \(\sup_n\|h_n\|<\infty\).
\end{examsubparts}
\end{examproblems}
\end{document}
